\chapter{Exercise 4A : Calculation of entropy of distributions}

The goal of this exercise is to calculate the entropy of the given discrete distributions, providing an analytical solution, if it exists, or a numerical one.

The distributions are the following

\begin{itemize}
	\item $ p(n) \propto e^{-\lambda n}  \qquad n = 0,1,2,3... $
	\item $ p(n) \propto e^{-n^2} \qquad n = 0,1,2,3...$
	\item $ p(n) \propto n^{-4} \qquad n = 1,2,3...$
	\item $ p(n) \propto \alpha^{n} \qquad n = 1,...,N $
	\item $ p(n) \propto (1 + n^2)^{-k} \qquad n = 0,1,2,... \qquad k = 1,2 $
\end{itemize}

\section{First distribution}

The first thing we do is normalizing the distribution to make sure that all the elements add up to 1. To do this we calculate the sum of the series analytically

\[ 
	\sum_{n=0}^{\infty} e^{- \lambda n}
\]

This is a geometric series with ratio equal to $ e^{- \lambda n} $.
This series only converge if the absolute value of the ratio is strictly less than 1, so the condition for convergence is
\[ 
	  |e^{- \lambda}| < 1   \Longrightarrow  \lambda > 0
\]

The sum of the distribution is calculated with the formula for calculating the sum of a geometric series, substituting the ratio $r$ with 
$e^{- \lambda}$

\[ 
	\sum_{n=0}^{\infty} e^{- \lambda n} =
	\frac{1}{1 - r} = \frac{1}{1 - e^{- \lambda}} =
	\frac{e^{\lambda}}{e^{\lambda} - 1}
\]
To normalize the distribution, we divide it for the inverse of the sum, so that the resuting series will add up to 1, so we obtain the following probability distribution
\[ 
	p(n) =
	\frac{e^{\lambda} - 1}{e^{\lambda}} \cdot e^{- \lambda n} =
	(1 - e^{-\lambda}) \cdot e^{- \lambda n}
\]
To calculate the entropy of the normalized distribution, we use the formula for calculating the entropy and apply basic logarithm properties to obtain
\[ 
	H = - \sum_{n=0}^{\infty} p(n) \cdot \log_2(p(n)) =
	- \sum_{n=0}^{\infty} 
		(1 - e^{-\lambda}) \cdot e^{- \lambda n}
	\cdot \log_2[
		(1 - e^{-\lambda}) \cdot e^{- \lambda n}
	] =
\]
\[ 
	- \frac{1 - e^{-\lambda}}{\ln(2)}  \cdot 
	\sum_{n=0}^{\infty}  e^{- \lambda n}
	\cdot 
		[ \ln(1 - e^{-\lambda}) - \lambda n]
	= 
\]
\[
	- \frac{1 - e^{-\lambda}}{\ln(2)}  \cdot 
	\left[ 
		\ln(1 - e^{-\lambda}) 
		\sum_{n=0}^{\infty} 
		e^{-\lambda n} -
		\lambda \sum_{n=0}^{\infty}
		n e^{-\lambda n}
	\right] =
\]
The first series has already been solved, the second one can be solved differentiating the geometric series in the following manner

\[ 
	\frac{d}{dr} \sum_{n=0}^{\infty} r^n =
	\frac{d}{dr} \left(
	\frac{1}{1 - r}
	\right)
	\rightarrow
	\sum_{n=0}^{\infty} nr^{n-1} =
	\frac{1}{(1 - r)^2}
	\rightarrow
	\sum_{n=0}^{\infty} nr^n =
	\frac{r}{(1 - r)^2}
\]
Substituting $e^{-\lambda}$ for $r$ in the solution we obtain
\[ 
	- \frac{1 - e^{-\lambda}}{\ln(2)}  \cdot 
	\left[ 
	\frac{\ln(1 - e^{-\lambda})}{1 - e^{-\lambda}} -
	\frac{\lambda e^{-\lambda}}{(1 - e^{-\lambda})^2}
	\right] =
	- \frac{\ln(1 - e^{-\lambda})}{\ln(2)}
	+ \frac{\lambda e^{-\lambda}}
	{\ln(2) (1 - e^{-\lambda}) } = 
\]
\[ 
	- \frac{\ln(e^{\lambda} - 1) - \lambda}
	{\ln(2)}
	+ \frac{\lambda}
	{\ln(2) (e^{\lambda} - 1) } =
	\frac{
		(1 - e^{\lambda})
		\left[
			\ln(e^{\lambda} - 1) - \lambda
		\right]
		+ \lambda
	}
	{\ln(2) \cdot (e^{\lambda} - 1)}
\]

\[
	H =
	\frac{
		(1 - e^{\lambda}) \cdot \ln(e^{\lambda} - 1) +
		\lambda e^{\lambda}
	}	
	{
		\ln(2) \cdot (e^{\lambda} - 1)
	}
\]

\subsection{Numerical examples}

To verify the results, we run numerical examples choosing some values for $\lambda$. \newline
First, we set $\lambda = 1$ and we obtain the following distribution
\[ 
	p(n) =
	\frac{e - 1}{e} \cdot e^{-n} =
	(1 - e^{-1}) \cdot e^{-n}
\]
Substituting $\lambda = 1$ in the analytical solution for the entropy we obtain
\[
	H = \frac{(1 - e) \cdot \ln(e - 1) + e}
	{\ln(2) \cdot (e - 1)} \approx 1.5013
\]
We compare this result to the one obtained using the formula for the entropy and adding the first 100 terms, and we see that it matches

\[
	H = - \sum_{n=0}^{\infty} p(n) \cdot \log_2(p(n)) =
	\sum_{n=0}^{\infty} (1 - e^{-1})\cdot e^{-n} \cdot
	\log_2((1 - e^{-1}) \cdot e^{-n}) \approx 1.5013
\]
Substituting $\lambda = 2$ we obtain the following distribution and entropy
\[
	p(n) = \frac{e^2 - 1}{e^2} \cdot e^{-2n} =
	(1 - e^{-2}) \cdot e^{-2n}
\]
\[
	H = \frac{(1 - e^2) \cdot \ln(e^2 - 1) + 2e^2}
	{\ln(2) \cdot (e^2 - 1)} \approx 0.6614
\]
We can see that the analytical solution is verified in this case too
\[
	H = - \sum_{n=0}^{\infty} p(n) \cdot \log_2(p(n)) =
	\sum_{n=0}^{\infty} (1 - e^{-2})\cdot e^{-2n} \cdot
	\log_2((1 - e^{-2}) \cdot e^{-2n}) \approx 0.6614
\]

\section{Second distribution}

For the second distribution, we again start with normalizing the distribution, but this time there is no analytical form for its sum, even if the series converge
\[
	\sum_{n=0}^{\infty} e^{-n^2} \approx 1.3863
\]
We use the numerical approximation as the normalizing factor and obtain
\[
	p(n) = \frac{1}{1.3863} \cdot e^{-n^2}
\]
The entropy of this distribution has no analytical form either, so we apply the formula and get a numerical approximation 

\[
	H = - \sum_{n=0}^{\infty} p(n) \cdot \log_2(p(n)) =
	- \frac{1}{1.3863} \sum_{n=0}^{\infty} e^{-n^2} \cdot 
	log2\left(\frac{1}{1.3863} \cdot e^{-n^2}\right) =
\]
\[
	- \frac{1}{1.3863 \cdot \ln(2)} \sum_{n=0}^{\infty} e^{-n^2} 
	 \left[ \ln \left(\frac{1}{1.3863}\right) - n^2 \right]
	 \approx 0.9315
\]

\section{Third distribution}

We calculate the sum of the distribution analytically, considering that the distribution is equal to the Riemann zeta function for $\zeta = 4$. There is a formula to calculate the value of the Riemann zeta function for even integers using the numbers of Bernoulli, so we get

\[ 
	\sum_{n=1}^{\infty} n^{-4} = \zeta(4) = 
	(-1)^{3} \cdot \frac{B_4(2\pi)^4}{2 \cdot 4!} = 
	\frac{\pi^4}{90} \approx 1.0823
\]
The normalized distribution is then

\[ 
	p(n) = \frac{90}{\pi^4 \cdot n^4}
\]
There is no closed form for the entropy of this distribution, so we calculate it numerically
\[
	H = - \sum_{n=1}^{\infty} \frac{90}{\pi^4 \cdot n^4} \cdot 
	\log_2\left(\frac{90}{\pi^4 \cdot n^4}\right) =
	- \frac{90}{\pi^4 \cdot \ln(2)} \cdot
	\sum_{n=1}^{\infty}
	\frac{\ln(90) - 4\ln(\pi) - 4\ln(n)} {n^4} = 
\]
\[ 
	- \frac{90}{\pi^4 \cdot \ln(2)} \cdot
	\left\{
	\left[ \ln(90) - 4\ln(\pi) \right] \zeta(4) -
	4 \sum_{n=1}^{\infty}
	\frac{\ln(n)}{n^4}
	\right\}
	\approx 0.4816
\]

\section{Fourth distribution}

For the fourth distribution, we consider the cases of the parameter $\alpha$ being equal or different to 1

\subsection{Case $\alpha = 1$}

If $\alpha = 1$ then the sum of the distribution is trivially $N$, so the normalized distribution is a vector of $N$ elements all equals to $1/N$. The entropy is then 
\[ 
	H = - \sum_{i=1}^{N} \frac{1}{N} \cdot \log_2\left(\frac{1}{N}\right)
	= - \log_2 \left(\frac{1}{N}\right) = \log_2(N)
\]
This value is also the maximum possible value for the entropy of a distribution of $N$ values

\subsection{Case $\alpha \ne 1$}

If $\alpha \ne 1$ then the sum of the distribution can be calculated using the formula for the sum of a finite geometric series with ratio $\alpha$, and substracting 1, since this series starts at the index 1, so it lacks the term $\alpha^0 = 1$

\[ 
	\sum_{n=1}^{N} \alpha^n = 
	\frac{\alpha^{N+1} - 1}{\alpha - 1} - 1 =
	\frac{\alpha \cdot (\alpha^N - 1)}{\alpha - 1}
\]
So the normalized distribution is

\[ 
	p(n) = \frac{\alpha - 1}{\alpha \cdot (\alpha^N - 1)} \cdot
		\sum_{n=1}^{N} \alpha^n
\]
The entropy is then
\[ 
	H = \sum_{n=1}^{N} p(n) \cdot \log_2(p(n)) =
\]

\section{Fifth distribution}

This distribution has two subcases, the one where $k = 1$ and the one where $k = 2$

\subsection{Case $k=1$}
In this case, the sum of the distribution has a closed form that can be expressed using hyperbolic functions

\[ 
	\sum_{n=0}^{\infty} \frac{1}{1 + n^2} = 
	\frac{1 + \pi \coth(\pi)}{2}
	\approx 2.0766
\]
The normalized distribution so is 
\[ 
	p(n) = \frac{2}{1 + \pi \coth(\pi)} \cdot
		\frac{1}{1 + n^2}
\]
The entropy of the normalized distribution doesen't have a closed form, so we have to calculate it numerically

\[ 
	K = \frac{2}{1 + \pi \coth(\pi)}
\]
\[ 
	H = - \sum_{n=0}^{\infty} p(n) \cdot \log_2(p(n)) =
	-  \sum_{n=0}^{\infty} \frac{K}{1 + n^2} \cdot
	\left[
	\log_2\left(\frac{K}{1 + n^2}\right) 
	\right] = 
\]

\[ 
	- K	\sum_{n=0}^{\infty} \frac{1}{1 + n^2} \cdot
	\left[
	\log_2(K)- \log_2(1 + n^2)
	\right] =
\]

\[ 
	- K \cdot \log_2(K) \cdot 
	\sum_{n=0}^{\infty} \frac{1}{1 + n^2}
	+ K \sum_{n=0}^{\infty}
	\frac{\log_2(1 + n^2)}{1 + n^2} =
\]
\[ 
	- K \cdot \log_2(K) \cdot 
	\frac{1}{K}
	+ K \sum_{n=0}^{\infty}
	\frac{\log_2(1 + n^2)}{1 + n^2} = 
	- \log_2(K)
	+ K \sum_{n=0}^{\infty}
	\frac{\log_2(1 + n^2)}{1 + n^2}
	\approx 2.5597
\]


\subsection{Case $k=2$}
In this case too, the sum of the distribution has a closed form that can be expressed using hyperbolic functions
\[ 
	\sum_{n=0}^{\infty} \frac{1}{(1 + n^2)^2} 
	= \frac{2 + \pi^2 \coth^2(\pi) + \pi \coth(\pi) - \pi^2}
	{4} \approx 1.3068
\]
The normalized distribution function is

\[ 
	p(n) = \frac {4} {2 + \pi^2 \coth^2(\pi) + \pi \coth(\pi) - \pi^2} \cdot
	\frac{1}{(1 + n^2)^2}
\]
The entropy of this distribution doesen't have a closed form, so it can only be calculated numerically

\[ 
	K = \frac {4} {2 + \pi^2 \coth^2(\pi) + \pi \coth(\pi) - \pi^2}
\]

\[ 
	H = - \sum_{n=0}^{\infty} p(n) \cdot \log_2(p(n)) =
	- \sum_{n=0}^{\infty} 
	\frac{K}{(1 + n^2)^2} \cdot 
	\log_2 \left[
	\frac{K}{(1 + n^2)^2}
	\right] =
\]
\[ 
	- \sum_{n=0}^{\infty} 
	\frac{K}{(1 + n^2)^2} \cdot
	\left[
		\log_2(K) - \log_2(1 + n^2)^2
	\right] =
\]
\[ 
	- K \cdot \log_2(K) \sum_{n=0}^{\infty}
	\frac{1}{(1 + n^2)^2} +
	K \sum_{n=0}^{\infty}
	\frac{\log_2(1 + n^2)^2}{(1 + n^2)^2} =
\]
\[ 
	- K \cdot \log_2(K) \cdot \frac{1}{K} +
	K \sum_{n=0}^{\infty}
	\frac{\log_2(1 + n^2)^2}{(1 + n^2)^2} =
\]
\[ 
	- \log_2(K) +
	K \sum_{n=0}^{\infty}
	\frac{\log_2(1 + n^2)^2}{(1 + n^2)^2}
	\approx 1.011
\]

